% Part: incompleteness
% Chapter: introduction
% Section: decidability

\documentclass[../../../include/open-logic-section]{subfiles}

\begin{document}

\olfileid{inc}{int}{dec}

\olsection{فیصلہ نہ ہو سکنا تے نامکملیت}

گوڈل دے نامکملیت قضیاں دے ثبوت لئی
نحو نوں عددی روپ دینا ضروری اے۔
پر ایہدے بغیر وی، صرف ایہ فرض کر
کے کہ اک نظریہ سارے !!{decidable}
تعلقّاں دی !!{represents} کردا اے،
کجھ اہم نتیجے حاصل ہو سکدے نیں۔
ایہ قطری دلیل رُکّن دے مسئلے دا
فیصلہ ناممکن ہون والی دلیل ورگی اے۔

\begin{thm}
جے $\Gamma$ اک بے تضاد نظریہ اے
جیہڑا ہر !!{decidable} تعلق دی
!!{represents} کردا اے، تاں
$\Gamma$ !!{decidable} نہیں۔
\end{thm}

\begin{proof}
فرض کرو $\Gamma$ !!{decidable}
ہووے۔ اسی دکھاواں گے کہ جے
$\Gamma$ ہر !!{decidable}
تعلق دی !!{represents} کردا اے،
تاں اوہ متضاد ہووے گا۔

% PNB-CMP-052: restore the indexed formula in the two diagonal membership-test fragments.
!!^{decidable} خاصیتاں، یعنی
یک-مقامی تعلق، اک آزاد
متغیر والے !!{formula}s نال
ظاہر کیتیاں جاندیاں نیں۔ فرض
کرو $!A_0(x)$، $!A_1(x)$، \dots،
ایہو جیہے سارے !!{formula}s ایس ترتیب نال
اک قابلِ حساب گنتی وار فہرست وچ آندے نیں۔
ہن ایہ سیٹ $D \subseteq \Nat$
ویکھو:
\[
D = \Setabs{n}{\Gamma \Proves \lnot !A_n(\num{n})}
\]
سیٹ~$D$ !!{decidable} اے، کیوں جے
اسی $n \in D$ جانچن لئی پہلے
$!A_n(x)$ نکال سکدے آں تے
اوہدے توں $\lnot !A_n(\num{n})$
بنا سکدے آں۔ ٹرم~$\num{n}$
نال $x$ دی ہر آزاد آمد
بدلنی، جیہڑی $!A_n(x)$ وچ
اے، تے
$!A_n(\num{n})$ دے اگے
$\lnot$ لانا صاف مشینی کم اے۔
فرض مطابق $\Gamma$ !!{decidable}
اے؛ سو اسی جانچ سکدے آں کہ
$\lnot !A_n(\num{n}) \in \Gamma$
اے یا نہیں۔ جے اے تاں
$n \in D$، تے جے نہیں تاں
$n \notin D$۔ ایس طرح
$D$ وی !!{decidable} اے۔

% PNB-CMP-053: normalize four malformed verbal glossary suffixes in this unit.
کیوں جے $\Gamma$ ساریاں
!!{decidable} خاصیتاں دی
!!{represents} کردا اے،
اوہ~$D$ دی وی
!!{represents} کردا اے۔
$D$ دی !!{represents} نظریے~$\Gamma$ وچ کرن والے
سارے !!{formula}s
فہرست $!A_0(x)$، $!A_1(x)$،
\dots، وچ شامل نیں۔ سو ایہو جیہا
عدد~$d$ چنو کہ $!A_d(x)$،
$D$ دی، نظریے~$\Gamma$ وچ،
!!{represents}
کردا ہووے۔ جے
$d \notin D$، تاں
$!A_d(x)$ ولوں~$D$
دی !!{represents} توں
$\Gamma \Proves \lnot !A_d(\num{d})$
ملدا اے۔ پر ایس دا مطلب
ایہ اے کہ $d$،
سیٹ~$D$ دی تعریف والی شرط
پوری کردا اے؛ یعنی
$d \in D$۔ ایہ
$d \notin D$ دے خلاف اے۔
ایس لئی برہانِ خلف نال
$d \in D$۔

جدوں $d \in D$ اے، تاں
سیٹ~$D$ دی تعریف توں
$\Gamma \Proves \lnot
!A_d(\num{d})$۔ دوجی
طرف، کیوں جے $!A_d(x)$،
$D$ دی، نظریے~$\Gamma$ وچ،
!!{represents} کردا اے،
$\Gamma \Proves !A_d(\num{d})$
وی اے۔ ایس لئی
$\Gamma$ متضاد اے۔
\end{proof}

\begin{explain}
پچھلا قضیہ دکھاندا اے کہ کوئی
بے تضاد نظریہ جیہڑا سارے
!!{decidable} تعلقّاں دی
!!{represents} کردا اے،
!!{decidable} نہیں ہو سکدا۔
اسی دکھاواں گے کہ
$\Th{Q}$ واقعی سارے
!!{decidable} تعلقّاں دی
!!{represents} کردا اے۔
ایس لئی $\Th{Q}$
نوں شامل کرن والے سارے
نظریے، مثلاً $\Th{PA}$
تے $\Th{TA}$، وی
ایہناں تعلقّاں دی نمائندگی
کردے نیں، تے ایس واسطے
اوہ وی !!{decidable}
نہیں۔ (ایہ سارے نظریے
معیاری نمونے وچ سچ نیں،
ایس لئی بے تضاد نیں۔)

ایس نتیجے توں پہلے
نامکملیت قضیے دا اک کمزور
روپ وی ملدا اے۔ ہر نظریہ
جیہڑا !!{axiomatizable}
تے !!{complete} اے،
!!{decidable} وی اے۔
ایس لئی بے تضاد نظریے
جیہڑے !!{axiomatizable}
نیں تے ساریاں
!!{decidable} خاصیتاں
دی !!{represents} کردے
نیں، !!{complete} نہیں
ہو سکدے۔
\end{explain}

\begin{thm}
جے $\Gamma$ !!{axiomatizable}
تے !!{complete} اے، تاں
اوہ !!{decidable} اے۔
\end{thm}

\begin{proof}
ہر متضاد نظریہ
!!{decidable} اے، کیوں جے
ایہو جیہے نظریے وچ سارے
!!{sentence}s ہوندے نیں۔
ایس لئی سوال «کیہ
$!A \in \Gamma$ اے؟»
دا جواب ہمیشہ «ہاں»
ہوندا اے، یعنی فیصلہ
کیتا جا سکدا اے۔

ہن فرض کرو $\Gamma$
بے تضاد اے، تے
!!{axiomatizable} تے
!!{complete} وی اے۔
!!{axiomatizable} ہون
دی وجہ توں $\Gamma$
!!{computably enumerable}
اے۔ نظریے~$\Gamma$ دے مسلّماں توں
جیہڑے درست !!{derivation}s نکلدے نیں،
اوہناں دی فہرست اک قابلِ حساب فنکشن
دے سکدا اے۔
ہر درست
!!{derivation} توں اوہ
!!{sentence} نکال سکدے
آں جیہڑا اوہ
!!{derive}s کردا اے؛
سو اک قابلِ حساب فنکشن
$\Gamma$ دے سارے
قضیاں دی فہرست بنا
دیندا اے۔ اک
!!{sentence}،
نظریے~$\Gamma$
دا قضیہ اُدوں تے صرف
اُدوں اے جدوں $\lnot !A$
قضیہ نہیں؛ ایہ
$\Gamma$ دی بے
تضادی تے
!!{complete} ہون
توں نکلدا اے۔
ہن $!A \in \Gamma$
دا فیصلہ ایویں کرو:
$\Gamma$ دے سارے
قضیے فہرست وچ آندے
جاؤن۔ جدوں $!A$
فہرست وچ آوے،
پتہ لگ جاندا اے
کہ $\Gamma \Proves !A$۔
جدوں $\lnot
!A$ آوے،
پتہ لگ جاندا اے
کہ $\Gamma \Proves/ !A$۔
$\Gamma$ دے
!!{complete} ہون
دی وجہ توں ایہناں
دوناں وچوں اک صورت
ضرور آوے گی؛
سو طریقہ آخرکار
جواب دے دیندا اے۔
\end{proof}

\begin{cor}
\ollabel{cor:incompleteness}
جے $\Gamma$ بے تضاد
تے !!{axiomatizable}
اے، تے ہر
!!{decidable}
خاصیت دی
!!{represents}
کردا اے، تاں
اوہ !!{complete}
نہیں۔
\end{cor}

\begin{proof}
جے $\Gamma$
!!{complete} ہوندا،
پچھلے قضیے توں
!!{decidable} ہوندا،
کیوں جے اوہ
!!{axiomatizable}
تے بے تضاد اے۔
پر $\Gamma$
ہر !!{decidable}
خاصیت دی
!!{represents}
کردا اے؛ پہلے
قضیے توں اوہ
!!{decidable}
نہیں۔ ایہ تضاد اے۔
\end{proof}

\begin{prob}
دکھاؤ کہ
$\Th{TA} = \Setabs{!A}{\Sat{N}{!A}}$
!!{axiomatizable}
نہیں۔ تُسی
ایہ فرض کر
سکدے او کہ
$\Th{TA}$
ساریاں فیصلہ
کرن جوگ
خاصیتاں دی
نمائندگی کردا اے۔
\end{prob}

جدوں اسی ثابت
کر لواں گے
کہ مثلاً
$\Th{Q}$
ساریاں
!!{decidable}
خاصیتاں دی
!!{represents}
کردا اے،
تاں نتیجہ
دسّے گا کہ
$\Th{Q}$
لازماً نامکمل
اے۔ پر ایس
دلیل وچ کسے
آزاد !!{sentence}
دی مثال نہیں
ملدی؛ صرف
ایہ ثابت
ہوندا اے کہ
ایہو جیہا !!a{sentence}
موجود اے۔
ایہو جیہی مثال
بناؤن لئی
نحو نوں عددی
روپ دینا تے
گوڈل دے
اصل ثبوت
دا خیال
استعمال کرنا
پینا اے۔
تے ظاہر اے
کہ سانوں
پہلا دعویٰ وی
ثابت کرنا
اے: $\Th{Q}$
واقعی ساریاں
!!{decidable}
خاصیتاں دی
!!{represents}
کردا اے۔

ایہ وی یاد
رکھنا چاہیدا
اے کہ ہر
\emph{دلچسپ}
نظریہ نامکمل
یا ناقابلِ
فیصلہ نہیں۔
بہت نظریے
دلچسپ ریاضیاتی
حقائق بیان
کرن جوگے
نیں، پر گوڈل
دے نتیجے
دیاں مطلوبہ
شرطاں پوری
نہیں کردے۔
مثلاً
$\Th{Pres} = \Setabs{!A \in \Lang{L_{A^+}}}{\Sat{N}{!A}}$
اوہ سیٹ اے جیہدے وچ حساب دی
اوہ زبان جیہدے وچ~$\times$ نہیں،
دے اوہی !!{sentence}s نیں
جیہڑے معیاری نمونے وچ سچ نیں۔
ایہ نظریہ
کامل وی اے
تے فیصلہ کرن
جوگ وی۔
ایہنوں
پریسبورگر
حساب کہندے
نیں؛ ایہ
قدرتی عدداں
بارے اوہ
سارے سچ
ثابت کردا اے
جیہڑے صرف
$\Obj{0}$،
$\prime$ تے~$+$
نال بیان
ہو سکدے نیں۔

\end{document}
